\documentclass[10pt,a4paper]{amsart}
\usepackage[margin=19mm]{geometry}
\usepackage{amsmath,amssymb,mathtools}
\usepackage[scr=rsfs,cal=euler]{mathalfa}
\usepackage{xfrac}
\usepackage[unicode,hidelinks,bookmarks=false]{hyperref}
\newcommand{\ie}{i.e.\ }
\newcommand{\cf}{cf.\ }
\newcommand{\resp}{respectively\ }
\newcommand{\Subsection}{\subsection}
\providecommand{\nobreakdash}{\nobreak\hbox{-}}
\setlength{\emergencystretch}{2em}
\multlinegap0pt
\usepackage{bm}
\usepackage{bbm}
\usepackage{dsfont}
\usepackage{xspace}
\usepackage{cancel}
\usepackage{mathtools}
\usepackage[shortlabels]{enumitem}
\usepackage{amsthm}
\makeatletter
\@ifundefined{theorem}{\newtheorem{theorem}{Theorem}}{}
\@ifundefined{lemma}{\newtheorem{lemma}[theorem]{Lemma}}{}
\makeatother
\newcommand{\Pbb}{\mathbb{P}}
\newcommand{\R}{\mathbb{R}}
\newcommand{\Scal}{\mathcal{S}}
\title[Moderate deviation principles for the current and the tagged]{Moderate deviation principles for the current and the tagged particle in the WASEP}
\author{Xiaofeng Xue, Linjie Zhao}
\date{Source: arXiv:2506.03924v1 (CC BY 4.0)}
\begin{document}
\maketitle

\noindent\emph{Source and licence.} Moderate deviation principles for the current and the tagged particle in the WASEP. Version arXiv:2506.03924v1 (CC BY 4.0), distributed under the Creative Commons Attribution 4.0 International licence, as recorded in the arXiv licence metadata for this exact version and shown on the abstract page https://arxiv.org/abs/2506.03924v1. Full rights notice: licenses/arxiv-2506.03924-CC-BY-4.0.txt.\par\smallskip

\noindent\textit{Editorial scope.} Counted unit: the Lemma whose source label is \texttt{lemma contraction principle of current} in the article (source0.tex lines 737--743, source label \texttt{lemma contraction principle of current}), delivered together with the article's own complete original proof of it (source0.tex lines 834--901, one \texttt{proof} environment of 68 lines). No mathematics is written, repaired or re-derived: statement and proof are the article's own text line for line. The proof is detached from the statement in the source: the article states the result at lines 737--743 and prints its proof at lines 834--901, and the proof environment's own header names the statement, so the association is the article's own. Required context retained uncounted: a t s (lines 101--107); lemma finite rate (lines 470--496); lem: finiteD mdp current (lines 565--569); exp tight error (lines 632--634); lemma 5.1 (lines 696--708); cov converg (lines 718--720); lemma 5.3 (lines 747--781). Every definition, display and label of those lines is retained unchanged. Omitted from this package and not redistributed: 1--100, 108--469, 497--564, 570--631, 635--695, 709--717, 721--736, 744--746, 782--833, 902--1191. Nothing inside a retained excerpt is omitted or altered: every retained body is delivered exactly as the article's lines are written, so no omission receipt of any kind accompanies this package. Citations: the retained excerpts carry no citation command at all, so no bibliography entry of the article is transcribed and no reference is redistributed. Reference adaptation: none was needed and none was made; every reference inside the delivered unit resolves inside it, so no unresolved reference or citation is produced. Printed slips kept literal: no printed word, symbol, hypothesis or number of the counted statement or of its proof was altered. Layout and class: the article's own class is \texttt{amsart} with packages including graphicx, amsfonts, amscd, epsfig, amsmath, amssymb, enumerate, bm, color, mathrsfs, geometry, hyperref, lineno, setspace; the wrapper is the lane's pinned \texttt{amsart} editorial wrapper at 10pt on A4, which restates the theorem-like environments the retained text needs and adds the article's own macro definitions that the retained text needs (\texttt{\textbackslash Pbb}, \texttt{\textbackslash R}, \texttt{\textbackslash Scal}), with extra wrapper packages (bm, bbm, dsfont, xspace, cancel, mathtools, enumitem). The delivered numbering is the unit's own. Assets: no retained span contains a figure environment, a graphics-inclusion command, a PDF-page inclusion, a table or a TikZ picture; no image file is copied into this package, the package declares no asset dependency and no image file is redistributed with it. Licence: version arXiv:2506.03924v1 of this article is distributed under the Creative Commons Attribution 4.0 International licence, as recorded in the arXiv licence metadata for this exact version and shown on the abstract page https://arxiv.org/abs/2506.03924v1. A full rights notice is delivered at licenses/arxiv-2506.03924-CC-BY-4.0.txt. Characters: every retained excerpt is pure ASCII, so the delivered engine, XeLaTeX with its default Latin Modern text font, reports no missing character.
\begin{equation}\label{a t s}
    a (t,s) = \begin{cases}
        \chi(\rho)\alpha|1-2\rho|\min\{t,s\} \quad &\text{if }\; \beta < 1,\\
        \chi(\rho)\left(f(t)+f(s)-f(|t-s|)\right) \quad &\text{if }\; \beta = 1,\\
        \chi(\rho)(\sqrt{t}+\sqrt{s}-\sqrt{|t-s|}) / \sqrt{2\pi}  \quad &\text{if }\; \beta > 1,
    \end{cases}
\end{equation}

\begin{lemma}\label{lemma finite rate} 
$(1)$ If $\mu_0 \in \Scal^\prime(\R)$ satisfies that $\mathcal{Q}_{\rm ini}(\mu_0)<+\infty$, then there exists $\varphi \in L^2(\R)$ such that
\[
\mu_0 (H) = \langle \varphi, H\rangle, \; \forall H \in L^2 (\R), \quad \text{and} \quad \mathcal{Q}_{\rm ini}(\mu_0)=\frac{1}{2\chi(\rho)} \|\varphi\|_{L^2 (\R)}^2.
\]
$(2)$ If $\mu\in D ([0, T], \Scal^\prime(\R))$ satisfies that $\mathcal{Q}^\beta(\mu)<+\infty$, then $\mu\in C([0, T], L^2(\R))$. Moreover, we have the following characterizations of  $\mu$ and the rate function.

\begin{itemize}
	\item If  $\beta>1$, then there exists $G\in \mathbb{H}$ such that $\mu$ is the unique weak solution to the PDE
	\[\begin{cases}
	\partial_t \mu(t,u)=\frac{1}{2}\Delta\mu(t,u)-\Delta G(t,u), \text{~}t> 0, u\in \mathbb{R},\\
	\mu(0,u) = \varphi (u), \; u \in \R,
	\end{cases}\]
where $\varphi$ is identified in $(1)$. Moreover, $\mathcal{Q}^\beta_{\rm dyn}(\mu)=[G,G] / (2 \chi(\rho))$.
\item If  $\beta=1$, then there exists $G\in \mathbb{H}$ such that $\mu$ is the unique weak solution to the PDE
\[\begin{cases}
	\partial_t \mu(t,u)=\mathcal{P}^* \mu(t,u)-\Delta G(t,u), \text{~}t\geq 0, u\in \mathbb{R},\\
	\mu(0,u) = \varphi (u), \; u \in \R.
\end{cases}\]
Here, $\mathcal{P}^*$ is the adjoint of $\mathcal{P}$ in $L^2 (\R)$. Moreover, $\mathcal{Q}^\beta_{\rm dyn}(\mu)=[G,G] / (2 \chi(\rho))$.
\item If  $\beta< 1$, then 
\[
\mu(t,u)=\varphi( u-\alpha(1-2\rho)t), \text{~}t\geq 0, u\in \mathbb{R}.
\]
Moreover, $\mathcal{Q}^\beta_{\rm dyn}(\mu)=0$.
\end{itemize}
\end{lemma}

\begin{lemma}\label{lem: finiteD mdp current}
For any $m \geq 1$ and for any $0 \leq t_1 < t_2 < \ldots < t_m \leq T$, the sequence of random vectors $\{\bar{J}^n_{-1,0} (t_i)/a_n, 1 \leq i \leq m\}_{n \geq 1}$ satisfies the MDP with decay rate $a_n^2 / n$ and with rate function $\mathcal{J}^{\beta,m} \equiv \mathcal{J}^\beta_{\{t_i\}_{i=1}^m}$, where  for $\mathbf{r} =(r_1,\ldots, r_m)^T\in \mathbb{R}^m$,
\[\mathcal{J}^{\beta,m}  (\mathbf{r}) = \frac{1}{2} \mathbf{r}^T A^{-1} \mathbf{r}.\]
Here, $A = A_{\{t_i\}_{i=1}^m} = (a(t_i,t_j))_{1 \leq i,j \leq m}$ is the $m \times m$ matrix with $a(\cdot,\cdot)$ defined in \eqref{a t s}. 
\end{lemma}

 \begin{equation}\label{exp tight error}
\limsup_{n \rightarrow \infty} \frac{n}{a_n^2} \log \Pbb_{\nu_\rho} \Big( \sup_{0 \leq t \leq T} |\langle \mu^n_t,F_l\rangle | >  \varepsilon \Big) = -\infty,
\end{equation}

\begin{lemma}\label{lemma 5.1}
Let $G\in \Scal(\R)$ . For any $m \geq 1$,  any $0 \leq t_1 < t_2 < \ldots < t_m \leq T$ and any $\mathbf{r}=(r_1,\ldots, r_m)^T\in \mathbb{R}^m$,
\[		\inf\left\{\mathcal{Q}^\beta_{t_m} (\mu):~ \langle \mu_{t_i} - \mu_{0}, G \rangle=r_i\text{~for all~}1\leq i\leq m\right\}
=\sup_{\bm{\xi} \in \R^m}\left\{\bm{\xi}^T \mathbf{r}-\frac{1}{2}\bm{\xi}^T \Sigma \bm{\xi}\right\},\]
	where $\Sigma := \Sigma_{\{t_k\}_{k=1}^m}(G)$ is a $m\times m$ symmetric matrix such that
	\[
	\Sigma (i,j)={\rm Cov}\left(\mathcal{Y}_{t_i}(G)-\mathcal{Y}_{0}(G), \mathcal{Y}_{t_j}(G)-\mathcal{Y}_{0}(G)\right)
	\]
	for any $1\leq i,j\leq m$. In particular, when $\Sigma$ is invertible, 
	\[
\inf\left\{\mathcal{Q}^\beta_{t_m} (\mu):~ \langle \mu_{t_i} - \mu_{0}, G \rangle=r_i\text{~for all~}1\leq i\leq m\right\}=\frac{1}{2} \mathbf{r}^T\Sigma^{-1} \mathbf{r} .
	\]
\end{lemma}

\begin{equation}\label{cov converg}
\lim_{l\rightarrow+\infty}{\rm Cov}\left(\mathcal{Y}_t(\tilde{G}_l)-\mathcal{Y}_0(\tilde{G}_l), \mathcal{Y}_s(\tilde{G}_l)-\mathcal{Y}_0(\tilde{G}_l)\right) = a(t,s),
\end{equation}

\begin{lemma}\label{lemma 5.3}
	Assume $\mu\in D ([0, T], \Scal^\prime(\R))$ satisfies that $\mathcal{Q}_{T}^\beta (\mu)<+\infty$.   Let  $G = G(\mu)$ and $\varphi = \varphi (\mu)$ be identified  in Lemma \ref{lemma finite rate}.
	\begin{enumerate}[(i)]
		\item For $\beta<1$,
		\[
		\int_0^{+\infty}[\mu(t,u)-\varphi (u)]du=\int_{-\alpha(1-2\rho)t}^0 \varphi (u) du, \quad 0\leq t\leq T.
		\]
		\item For $\beta>1$,
		\begin{align*}
			\int_0^{+\infty}&[\mu(t,u)-\varphi (u)]du\\
			&=\int_{\R}\varphi (u) V_t(u)du-\int_0^t\left(\int_{\R}\partial_{uu}^2G (s,u)\mathbb{P}(B_{t-s}+u\geq 0)du\right)ds, \quad 0 \leq t \leq T,
		\end{align*}
		where $\{B_t\}_{t\geq 0}$ is the one dimensional standard Brownian motion starting from the origin and 
		\[
		V_t(u)=\mathbb{P}(B_t+u\geq 0)-\chi_{\{u\geq 0\}}=
		\begin{cases}
			\mathbb{P}(B_t\geq |u|) & \text{~if~}u<0,\\
			-\mathbb{P}(B_t\geq |u|) & \text{~if~}u\geq 0. 
		\end{cases}
		\]
		\item For $\beta=1$,
		\begin{align*}
			&\int_0^{+\infty} [\mu(t,u)-\varphi (u)] du=\int_{\R}\varphi (u)R_t(u)du\\
			&-\int_0^t\left(\int_{\R}\partial_{uu}^2 G (s,u)\mathbb{P}(B_{t-s}+u+\alpha(t-s)(1-2\rho)\geq 0)du\right)ds, \quad 0 \leq t \leq T,
		\end{align*}
		where 
		\[
		R_t(u)=\mathbb{P}(B_t\geq -u-\alpha(1-2\rho)t)-\chi_{\{u\geq 0\}}=
		\begin{cases}
			\mathbb{P}(B_t\geq -u-\alpha(1-2\rho)t) & \text{~if~}u<0,\\
			-\mathbb{P}(B_t\geq u+\alpha(1-2\rho)t) & \text{~if~}u\geq 0. 
		\end{cases}
		\] 
	\end{enumerate}
\end{lemma}

\begin{lemma}\label{lemma contraction principle of current}
For any $\beta\geq 0$ and any $\mathbf{r}=(r_1,\ldots, r_m)^T\in \mathbb{R}^m$, 
\[
\inf\left\{\mathcal{Q}^\beta_{t_m} (\mu):~\int_0^{+\infty} [\mu(t_i,u)-\mu(0, u)] du=r_i\text{~for all~}1\leq i\leq m\right\}\geq \frac{1}{2} \mathbf{r}^T A^{-1} \mathbf{r},
\]
where the matrix $A$ was defined in Lemma \ref{lem: finiteD mdp current}.
\end{lemma}

\begin{proof}[Proof of Lemma \ref{lemma contraction principle of current}]
	We only deal with the  case $\beta = 1$ since the remaining cases follow from the same analysis.   It suffices to show that, if  $\mathcal{Q}^\beta_{t_m}(\mu)<+\infty$ and $\int_0^{+\infty} [\mu(t_i,u)-\mu(0, u) ] du=r_i$ for all $1\leq i\leq m$, then 
	\[
	\mathcal{Q}^\beta_{t_m} (\mu)\geq \frac{1}{2} \mathbf{r}^T A^{-1} \mathbf{r}.
	\]
	
	We first prove the above inequality for $\mu$ such that
	\begin{equation}\label{equ condition 5.1}
		\varphi \in C^\infty_c(\R) \text{~and~} G \in C_c^{1,+\infty}([0, T]\times \R),
	\end{equation}
where $\varphi$ and $G$ were identified in Lemma \ref{lemma finite rate}. Let $\tilde{G}_l\in \Scal(\R)$ be the approximation of $G_l$ introduced before \eqref{exp tight error}. By Lemma \ref{lemma 5.1},
	\[
	\mathcal{Q}^\beta_{t_m} (\mu)\geq \frac{1}{2} (\mathbf{r}^l)^T\left(\Sigma_{\{t_k\}_{k=1}^m}(\tilde{G}_l)\right)^{-1}\mathbf{r}^l,
	\]
	where $\mathbf{r}^l=(r_1^l,\ldots, r_m^l)^T \in \R^m$ with $r_i^l=\int_{\R}(\mu(t_i, u)-\varphi (u))\tilde{G}_l(u)du$ for all $1\leq i\leq m$.
	We claim that
	\begin{equation}\label{equ claim proof of Lemma 5.2}
		\lim_{l\rightarrow+\infty} r_i^l= r_i.
	\end{equation}
Then, by \eqref{cov converg},
	\[
	\mathcal{Q}_{t_m} (\mu)\geq \limsup_{l\rightarrow+\infty} \frac{(\mathbf{r}^l)^T\left(\Sigma_{\{t_k\}_{k=1}^m}(\tilde{G}_l)\right)^{-1}\mathbf{r}^l}{2} = \frac{1}{2} \mathbf{r}^T A^{-1} \mathbf{r}.
	\]

	Now we prove \eqref{equ claim proof of Lemma 5.2}.  
	As in the proof of Lemma \ref{lemma 5.3}, for any $H\in L^2(\R)$, 
	\begin{align}\label{equ 5.3}
		\int_{\R}\left(\mu(t,u)-\varphi (u)\right)H(u)du=&\int_{\R}\varphi (u)\left(\mathbb{E}[H(B_t+u+\alpha(1-2\rho)t)]-H(u)\right)du \\
		&-\int_0^t\left(\int_{\R}\partial_{uu}^2G (s,u)\mathbb{E}[H\left(B_{t-s}+\alpha(1-2\rho)(t-s)+u\right)] du\right)ds\notag. 
	\end{align}
	By Equation \eqref{equ 5.3}, Cauchy-Schwarz inequality and the fact that $G \in C_c^{1,+\infty}([0, T]\times \R)$, 
	\[
	\lim_{l\rightarrow+\infty}\int_{\R}(\mu(t, u)- \varphi (u))(\hat{G}_l(u)-G_l(u))du=0. 
	\]
	Hence, to check Equation \eqref{equ claim proof of Lemma 5.2}, we only need to show that
	\begin{equation}\label{equ 5.4}
		\lim_{l\rightarrow+\infty}\int_{\R}(\mu(t, u)-\varphi (u)) G_l (u)du=\int_0^{+\infty}\mu(t, u)-\varphi (u) du.
	\end{equation}
	By Lemma \ref{lemma 5.3} and Equation \eqref{equ 5.3},
	\begin{align}\label{equ 5.5}
		&\int_{\R}(\mu(t, u)-\mu(0,u))G_l(u)du-\int_0^{+\infty}\mu(t, u)-\mu(0,u)du \notag\\
		&=\int_{\R} \varphi (u) F_l(u)du-\int_0^t\left(\int_{\R}\partial_{uu}^2 G (s,u)Q_l(s,u)du\right)ds,
	\end{align}
	where
	\begin{align*}
		F_l(u) &=\mathbb{E}[G_l(B_t+u+\alpha(1-2\rho)t)]-G_l(u)-R_t(u),\\
		Q_l(s,u)&=\mathbb{E}[G_l\left(B_{t-s}+\alpha(1-2\rho)(t-s)+u\right)]-\mathbb{P}(B_{t-s}+u+\alpha(t-s)(1-2\rho)\geq 0).
	\end{align*}
	According to the definition of $G_l$, we have
	\begin{align*}
		|Q_l(s,u)|&\leq \frac{\mathbb{E}|B_{t-s}+u+\alpha(1-2\rho)(t-s)|}{l}+\mathbb{P}\left(B_{t-s}+u+\alpha(t-s)(1-2\rho)\geq l\right).  
	\end{align*}
	Therefore, $\lim_{l\rightarrow+\infty}Q_l (s,u)=0$ uniformly on any compact set in $ [0, t]\times \R$. Similarly, $\lim_{l\rightarrow+\infty}F_l=0$ uniformly on any compact set in $\R$. Then, Equation \eqref{equ 5.4} holds according to Equation \eqref{equ 5.5} and the fact that $\varphi$ and $G$ have compact support. 
	
	For $\mu$ not satisfying \eqref{equ condition 5.1}, the idea is to approximate $\varphi (\mu)$ and $G (\mu)$ by $\varphi^\varepsilon \in C^\infty_c (\R)$ and $G^{\varepsilon} \in C_c^{1,\infty} ([0,T] \times \R)$ respectively.  Precisely speaking, let  $\varphi^\varepsilon  \rightarrow \varphi$ in $L^2(\R)$ and $G^{\varepsilon}  \rightarrow G$ in $\mathbb{H}$ as $\varepsilon \rightarrow 0$.  Let $\mu^\varepsilon \in D ([0, t_m], \Scal^\prime(\R))$ be given in Lemma \ref{lemma finite rate} associated with $\varphi^\varepsilon $ and $G^{\varepsilon}$. 
	Then, by Lemma \ref{lemma finite rate},
	\begin{align}\label{equ limit in proof of Lemma 5.2}
		\lim_{\varepsilon\rightarrow 0}\mathcal{Q}^\beta_{t_m}(\mu^\varepsilon)&=  \lim_{\varepsilon\rightarrow 0} \frac{1}{2\chi (\rho)} \left(\|\varphi^\varepsilon\|_{L^2(\R)}^2+ [G^{\varepsilon},G^{\varepsilon}]\right) \notag\\
		&=\frac{1}{2\chi (\rho)} \left(\|\varphi\|_{L^2(\R)}^2+ [G,G]\right)=\mathcal{Q}_{t_m}(\mu)
	\end{align}
	Since $\mu^\varepsilon$ satisfies \eqref{equ condition 5.1},  we have shown that
	\[
	\mathcal{Q}^\beta_{t_m} (\mu^\varepsilon)\geq \frac{1}{2} (\mathbf{r}^\varepsilon)^T A^{-1} \mathbf{r}^\varepsilon,
	\]
	where $\mathbf{r}^\varepsilon=(r^\varepsilon_1,\ldots, r^\varepsilon_m)^T$ with $r^\varepsilon_i=\int_0^{+\infty}\mu^\varepsilon(t_i,u)-\varphi^\varepsilon (u) du$ for all $1\leq i\leq m$.  Thus,
	\[\mathcal{Q}_{t_m}^\beta (\mu) \geq    \limsup_{\varepsilon\rightarrow 0} \frac{1}{2} (\mathbf{r}^\varepsilon)^T A^{-1} \mathbf{r}^\varepsilon.\]
By Lemma \ref{lemma 5.3},  $\lim_{\varepsilon \rightarrow 0} \mathbf{r}^\varepsilon = \mathbf{r}$, thus concluding the proof. 
\end{proof}

\end{document}
